% Lead author: Reza Iraji
% Paper coauthors: Felipe Galarza-Jimenez and Majid Zamani

% Insert immediately before Theorem 1.
Although path-completeness is a property of the labeled graph alone, the existence of a PC-CC is what couples that graph to the switched dynamics. The theorem below characterizes when the local graph-indexed certificate conditions are sound for arbitrary switching.

% Optional lemma to shorten the sufficiency proof.
\begin{lemma}[Path lifting]\label{lem:path-lifting}
Let $\mathcal G=(\mathcal V,\mathcal E,\Sigma)$ be finite and path-complete. For every infinite word $\boldsymbol\sigma\in\Sigma^\omega$, there exists an infinite node sequence $(v_t)_{t\in\mathbb N}$ such that $(v_t,\boldsymbol\sigma(t),v_{t+1})\in\mathcal E$ for all $t$. The sequence need not be unique.
\end{lemma}
\begin{proof}
At depth $T$, collect all graph paths labeled by the length-$T$ prefix of $\boldsymbol\sigma$. Path-completeness makes every level nonempty, and finiteness of $\mathcal V$ makes the prefix tree finitely branching. K\"onig's infinity lemma yields an infinite branch.
\end{proof}

% Keep the original PC-CC theorem and direct SOS proposition unchanged in meaning.
\begin{proposition}[SOS sufficient conditions for PC-CCs]\label{prop:sos-pccc}
The SOS conditions stated below imply PC-CC1--PC-CC3 on the corresponding compact semialgebraic domains.
\end{proposition}

% New minimal extension used only for the four-visit platoon case.
\begin{proposition}[Target-local ranked path-complete certificate]\label{prop:target-local-ranked}
Let $\mathcal G=(\mathcal V,\mathcal E,\Sigma)$ be finite and path-complete, and let $\mathcal T\subseteq\mathcal X_{\mathrm{vf}}$ contain every state in $\mathcal X_{\mathrm{vf}}$ reachable from $\mathcal X_0$. Suppose there are functions $C_{p,q}:\mathcal X\times\mathcal X\to\mathbb R$ and a rank $R:\mathcal T\to\mathbb R_{\ge0}$ such that
\begin{align*}
C_{p,q}(x,f_\sigma(x))&\ge0 &&\forall (p,\sigma,q)\in\mathcal E,\ x\in\mathcal X,\\
C_{q,r}(f_\sigma(x),y)\ge0&\Longrightarrow C_{p,r}(x,y)\ge0
&&\forall (p,\sigma,q)\in\mathcal E,\ r\in\mathcal V,\ x\in\mathcal X,\ y\in\mathcal T,\\
C_{q,q}(y,y')\ge0&\Longrightarrow R(y')\le R(y)-1
&&\forall q\in\mathcal V,\ y,y'\in\mathcal T.
\end{align*}
Then every trajectory starting in $\mathcal X_0$ visits $\mathcal X_{\mathrm{vf}}$ only finitely many times.
\end{proposition}
\begin{proof}
Fix any switching word and use Lemma~\ref{lem:path-lifting} to select a compatible infinite node path. By repeated use of the second implication, starting from the first inequality on the last edge of any finite path segment, every pair of visit states $x_{t_i},x_{t_j}\in\mathcal T$ with $t_i<t_j$ satisfies $C_{v_{t_i},v_{t_j}}(x_{t_i},x_{t_j})\ge0$. If there were infinitely many visits, finiteness of $\mathcal V$ would give a node $q$ occurring at infinitely many visit times. Applying the third implication to successive visits on this repeated-node subsequence decreases the nonnegative rank by at least one each time, a contradiction.
\end{proof}

% SOS realization of Proposition~\ref{prop:target-local-ranked}.
For polynomial $C_{p,q}$ and $R$ on compact semialgebraic domains, each implication in Proposition~\ref{prop:target-local-ranked} is enforced by a Positivstellensatz/SOS certificate. In particular, the premise polynomial is retained as a domain generator rather than canceled by fixing a unit multiplier. We refer to the recurrent-node rank implication as SOS-R.

For a graph with $|\mathcal V|$ nodes and $|\mathcal E|$ edges, the direct original-PC-CC encoding contains $|\mathcal E|$ SOS-1 conditions, $|\mathcal E||\mathcal V|$ SOS-2 conditions, and $|\mathcal V|^3$ SOS-3 conditions before domain-multiplier expansion. The target-local ranked formulation used in the platoon instead has $|\mathcal E|$ one-step closure conditions, $|\mathcal E||\mathcal V|$ target-local propagation conditions, and $|\mathcal V|$ recurrent-node rank conditions.

% Add to the limitations discussion.
The ranked formulation is a separate sufficient condition and is not identified with PC-CC3. The scalable graph-indexed benchmark in the numerical section independently exercises the original PC-CC3 condition non-vacuously. We also do not claim completeness within any prescribed polynomial degree or SOS relaxation order.
